% !TEX root = ../algebra_02.tex

\section{Смежные классы, нормальные подгруппы.}

\subsection{Порожденная подгруппа.}

\begin{Definition}{Операции с подгруппами}{}
    
    $G$~--- группа.
    
    $P, Q \subset G$:
    \[
        P\cdot Q := \{p\cdot q\ \ | \ \ p \in P, q \in Q \}        
    \]
    $p \in P$:
    \[
        p\cdot Q := \{p\cdot q \ \ |\ \  q \in Q \}
    \]
    \[
        Q\cdot p := \{ q\cdot p\ \ | \ \ q \in Q \}
    \]

    $P \subset G$
    \[
        P^{-1} := \{ p^{-1} \ \ | \ \ p \in P \}
    \]
\end{Definition}

\begin{Definition}{Порожденная подгруппа}{}

    $M \subset G, G$~--- группа.
    \textbf{Подгруппа, порожденная} $M$~--- $H_0 < G$, т.ч.:

    \begin{enumerate}
        \item $M \subset H_0$
        \item $H < G, M \subset H \Rightarrow H_0 \subset H$
    \end{enumerate}

    $<M> :=$ подгруппа, порожденная $M$.
    
\end{Definition}

\begin{remark}

    Если $H_0, H_1$ удовлетворяют условиям 1 и 2, то 
    \[ H_0 \subset H_1, H_1 \subset H_0 \Leftrightarrow H_0 = H_1 \]
\end{remark}

\begin{Proposition}{}{}
    $G$~--- группа. $M < G \Rightarrow$
    \[ \bigcap_{\substack{ H < G \\ M \subset H}} H =: <M> \]
\end{Proposition}

\begin{proof}
    \[ H_0 := \bigcap_{\substack{ H < G \\ M \subset H}} H \]
    \begin{enumerate}
        \item $H_0 < G$
        
        (пересечение любого количества подгрупп - подгруппа)
        \item $M \subset H_0$
        
        (т.к. $M \subset H \ \ \forall \ \ H$)

        \item $\forall H < G, M \subset H \Rightarrow H \supset H_0$
        
        (т.к. $H_0$~--- пересечение всех таких H)
    \end{enumerate}   
\end{proof}

\begin{Proposition}{}{}
    \[
    <\{g_1, g_2, \ldots, g_n \} > =: <g_1, g_2, \ldots, g_n>
    \]
\end{Proposition}

\begin{remark}
    Легко видеть, что $<g_0>$~--- подгруппа порожденная $\{g_0\}$.
\end{remark}

\begin{Proposition}{}{}

    $G$~--- группа. $M < G$, тогда:
    \[
    <M> = \{ g_1\cdot g_2\textsc{} \cdot \ldots \cdot g_n \ \ | \ \ n \ge 0; \ \ g_i \in M \cup M^{-1} \} 
    \]

    \begin{remark}
        $(n = 0 \Rightarrow g_1 \cdot \ldots \cdot g_n = e)$
    \end{remark}
\end{Proposition}

\begin{proof}
    \[
    H_0 := <M> = \{ g_1\cdot g_2\textsc{} \cdot \ldots \cdot g_n \ \ | \ \ n \ge 0; \ \ g_i \in M \cup M^{-1} \} 
    \]

    \begin{itemize}
        \item Очевидно $M \subset H_0$
        \item Проверим $H_0 < G$:

        \begin{enumerate}
            \item $e \in H_0$ (в силу замечания)
            \item $g, g' \in H_0 \Rightarrow$
            \[ g = g_1\cdot \ldots \cdot g_n \in M \cup M^{-1}; \qquad g' = g'_1\cdot \ldots g'_m \in M \cup M^{-1} \]
            \[ g\cdot g' = g_1\cdot \ldots \cdot g_n \cdot g'_1 \cdot \ldots \cdot g'_m \in M \cup M^{-1} \]
            \item $g \in H_0 \Rightarrow$
            \[ g = g_1 \cdot \ldots \cdot g_n, \ \ g_i \in M \cup M^{-1} \Rightarrow \]
            \[ g^{-1} = g_n^{-1} \cdot \ldots \cdot g^{-1}_1, g^{-1}_i \in M^{-1} \cup M = M \cup M^{-1} \Rightarrow \]
            $g^{-1} \in H_0$
        \end{enumerate}

        \item $H < G, M \subset H \Rightarrow H \supset H_0$:

        \[ H \supset M \Rightarrow H \supset M^{-1} \Rightarrow H \ni g_1\cdot \ldots \cdot g_n \qquad g_i \in M \cup M^{-1} \Rightarrow H \supset H_0 \]
    \end{itemize}
\end{proof}

\begin{remark}

    Так как в общем случае элементы не коммутируют, подгруппы, порожденные из двух элементов уже могут давать довольно сложные комбинации:

    $<s, t> \ni s^5\cdot t^{-2} \cdot s \cdot t^3 \cdot s^{-1}$
\end{remark}

\begin{Example}{Группа симметрий правильного n-угольника}{}

    Группа $D_n$ порождена двумя элементами: $S_n = <R_n, S_n>$ (поворот на $k\cdot \frac{2\pi}{n}$ радиан и осевая симметрия)

    \[   
    D_n = \{ R_i, S_i \ \ | \ \ 0 \le i \le n-1 \}
    \]

    $R_k = R^k_1 \Rightarrow R_i \circ R_j = R_{i+j}$

    Из геометрии можно увидеть: $S_i \circ S_j = R_{i-j}$

    Построив Латинский квадрат, получим, что $R_i \circ S_j = S_x$ (т.к. $R_i \circ R_j$ = $R_{i+j}$)
    
    \[  
    R_i \circ S_j = S_x \Rightarrow R_{-i} \circ R_i \circ S_j \circ S_x = R_{-i} \circ S_x \circ S_x
    \]
    \[
    R_{j-x} = S_j \circ S_x = R_{-i} \Rightarrow -i = j - x \Rightarrow x = i +j
    \]

    $R_i \circ S_j = S_{i+j}$

    $(S_i \circ R_j)^{-1} = R^{-1}_j \circ S^{-1}_i = R_{-j} \circ S_i = S_{i-j}$
    
    \[ H_0 = <R_1, S_0> \qquad \forall \ \ k \st R^k_1 \in H_0 \qquad R^k_1\circ S_0 = S_k \in H_0 \Rightarrow H_0 = D_n \]
\end{Example}

\begin{Example}{Перестановки}{}

    $H_0 = <\sigma_1, \tau>, \qquad \tau = (12)$

    $\sigma \tau \sigma^{-1} = (23)\Rightarrow$

    \[ \tau' = (k, k+1) = \sigma^{k-1}\tau \sigma^{1-k} \in H_0 \qquad \forall \ \ k \Rightarrow \]
    $H_0 = <\sigma, \tau> = S_n$
\end{Example}

\subsection{Смежные классы.}

Мотивация~--- хотим научиться понимать, насколько один элемент "отличается" от другого.

Для двух элементов $g', g \in G$, если $g' = g \cdot h$ для некоторого $h \in H < G$ - $g'$ "отличается" от $g$ на $h$.

Теперь формально:

$H < G$. Введем на $G$ отношение $\sim$:

\[
g' \sim g \Leftrightarrow g' = gh, \qquad h \in H
\]

\begin{Statement}{$\sim$~--- отношение эквивалентности}{}
    
    Проверим:

    \begin{enumerate}
        \item $g \sim g \st g = g\cdot e, \qquad e \in H$ (т.к. $H < G$)
        \item $g' \sim g \Rightarrow g \sim g' \st g' = gh \Rightarrow g = g'h^{-1}$ (т.к. $h\in H \Rightarrow h^{-1} \in H$)
        \item $g_1 \sim g_2, \ \ g_2 \sim g_3 \Rightarrow g_1 \sim g_3 \st$
        \[ g_1 = g_2h, \ \ g_2 = g_3h', \qquad h, h' \in H\]
        \[ g_1 = g_3\cdot h\cdot h', \qquad h\cdot h' = \tilde{h} \in H\]
        \[ g_1 = g_3 \tilde{h} \Leftrightarrow g_1 \sim g_3 \]
    \end{enumerate}
\end{Statement}

Рассмотрим множество элементов, эквивалентных $g$: 

\begin{Definition}{Класс эквивалентности}{}

    Класс эквивалентности $g$ (на $G$) = $[g]$ = $\{ gh \ \ | \ \ h \in H \} := gH$
\end{Definition}

В общем случае $gh \not = hg$, поэтому в нашем случае:

$\sim$~--- \textbf{отношение левой смежности по подгруппе} $H$

\begin{Definition}{Смежный класс}{}

    $gH$~--- \textbf{левый смежный класс} $g$ по $H$, или же \textbf{левый класс смежности}.
\end{Definition}

\begin{Definition}{Множжество классов смежности}{}

    $G/H = G/\sim := \{ gH \ \ | \ \ g \in G \}$~--- множество всех левых смежных классов по $H$.
\end{Definition}

\begin{Definition}{Индекс подгруппы}{}
    
    Количество левых классов смежности в группе:

    $|G/H| =: (G:H)$~--- \textbf{индекс подгруппы $H$ в группе $G$}.
\end{Definition}

\begin{Example}{}{}

    $\ZZ/_{n\ZZ}$~--- класс вычетов по модулю n.
    
    Индекс = $(\ZZ : n\ZZ) = n$.
\end{Example}

\begin{Definition}{Правый класс смежности}{}

    Правый класс смежности $g$ по $H =: Hg$

    Множество правых смежных классов $= \{ Hg \ \ | \ \ g \in G\} =: H \backslash G$
\end{Definition}

(Стоит обращать внимание на контекст, чтобы различать операцию разности множеств $H \setminus G$ и множество правых классов смежности $H \backslash G$)

\begin{remark}
    Разные авторы могут определять левые и правые классы смежности наоборот.
\end{remark}

\begin{Proposition}{}{}
    $|H \backslash G| = (G:H) = |G \slash H|$
\end{Proposition}

\begin{proof}

    $\alpha: \underset{p \rightarrow p^{-1}}{G \slash H \longrightarrow H \backslash G}$

    Проверим $p \in G \slash H \Rightarrow p^{-1} \in H \backslash G \st$

    \[
     p = gH \Rightarrow p^{-1} = H^{-1}g^{-1} = Hg^{-1} \in H \backslash G \ \ (g^{-1} \in G)
    \]

    Аналогично построим корректное отображение

    $\beta: \underset{q^{-1} \leftarrow q}{G \slash H \longleftarrow H \backslash G}$

    $\beta \circ \alpha =: id_{G \slash H}, \qquad \alpha \circ \beta = id_{H \backslash G}$

    Существует обратное отображение $\beta, \Rightarrow \alpha$~--- биекция, а значит множества равномощны.  
\end{proof}

\textbf{Следствие:} выводы для левой группы смежности будут работать и для правой.

\begin{Theorem}{Теорема Лагранжа}{}

    $|G| < \infty, \qquad H < G \Rightarrow$
    \[ |G| = |H|\cdot (G:H) \]
\end{Theorem}
\begin{proof}
    \[ |G| = \sum_{P \in G \slash H } |P| \]
    $P = gH$, построим отображение $\eps \st H \longrightarrow gH, h \rightarrow gh$

    $\eps$~--- сюрьективно по определению $gH$. При этом $gh_1 = gh_2 \overset{\text{сокращение в группе}}{\Rightarrow} h_1 = h_2$

    $\Rightarrow \eps$~--- биекция. Тогда $|gH| = |H|$

    \[ |G| = \sum_{P \in G \slash H } |P| = \sum_{P ]in G \slash H} |H| = |H| \cdot |G \slash H| \]
\end{proof}

\textbf{Следствия:}

\begin{enumerate}
    \item $H < G$, \ \ $|G| < \infty \Rightarrow |H| | |G|$
    \item $g \in G, \ \ |G| < \infty \Rightarrow (ord(g) = |<g>| = ord(H) = |H|) | |G|$
    \item $g \in G, \ \ |G| < \infty \Rightarrow g^{|G|} = e$
    \item $P \ni p = |G| \Rightarrow G$~--- цикл.
    \[ g \in G \setminus \{e\}, \qquad ord(g) | |G| \Rightarrow ord(g) | p \]
    \[ g \not = e \Rightarrow ord(g) \not = 1, ord(g) = P \Rightarrow ord(g) = |G| \Rightarrow\]
    $G$~--- цикл и порождается каждым своим элементом.
\end{enumerate}

\begin{Example}{Применение к теории чисел}{}
    Используем следствие 3:

    $G = (\ZZ \slash n\ZZ)^*$

    $|G| = \vphi(n)$~--- количество обратимых классов

    По следствию: $(\overline{a})^{\vphi(n)} = \overline{1} = e$
    \[\overline{a^{\vphi(n)}} = \overline{1} \Rightarrow a^{\vphi(n)} = 1 \mod n \qquad (\gcd(a, n) = 1) \]~--- теорема Эйлера.
\end{Example}

\subsection{Нормальные подгруппы.}

Мотивация:

Получить свойство $(g_1H)\cdot (g_2H) = g_1g_2H$

\[ (g_1H)\cdot (g_2H) = g_1g_2H \]
\[ g_1(Hg_2)H = g_1g_2H \]
\[ (Hg_2)H = g_2H \Leftrightarrow Hg_2 = g_2H \]

~--- полезно получить $Hg_2 = g_2H$

\begin{Proposition}{}{}

    $H$~--- подгруппа $G$. Тогда эквивалентны:

    \begin{enumerate}
        \item $\forall \ \ g \in G, \ \ \forall \ \ h \in H \st gh = h'g$ для некоторого $h' \in H$
        \item $\forall \ \ g \in G, \ \ \forall h \in  H \st ghg^{-1} \in H$~--- сопряженный к $h$ при помощи $g$ содержится в $H$.
        \item $\forall g \in G \st Hg = gH$
        \item $\forall \ \ g \in G, \st gHg^{-1} \subset H$
        \item $\forall \ \ g \in G, \st gHg^{-1} = H$
    \end{enumerate}
\end{Proposition}

\begin{proof}
    
    $1 \rightarrow 2$:

    \[ gh = h'g, \ \ h' \in H \Rightarrow ghg^{-1} = h' \in H \]
    $2 \rightarrow 4$:

    \[ gHg^{-1} = \{ ghg^{-1}, h \in H \} \] (по введеной операции произведения элемента на множество)

    \[ ghg^{-1} \in H \Rightarrow gHg^{-1} \subset H \]

    $4 \rightarrow 5$:

    \[ gHg^{-1} < H; \qquad g' := g^{-1} \Rightarrow g'^{-1}Hg' \subset H \Rightarrow \]
    \[ H \subset g'Hg'^{-1} \Rightarrow gHg^{-1} = H \]

    $5 \rightarrow 3$:

    \[ gHg^{-1} = H \Rightarrow gHg^{-1}\cdot g = Hg \Rightarrow gH = Hg \]

    $3 \rightarrow 1$:

    \[ gh \in gH = Hg \Rightarrow gh \in Hg \Rightarrow gh = h'g, \qquad h' \in H \]

\end{proof}

\begin{Definition}{Нормальная подгруппа}{}
    Подгруппу $H < G$ с свойства 1-5 назовем \textbf{нормальной подгруппой}.

    $H \vartriangleleft G$
\end{Definition}

\begin{Example}{}{}
    \begin{enumerate}
        \item $\{e\} \vartriangleleft G \qquad (g\{e\} = \{g\} = \{e\}g)$
        \item $G \vartriangleleft G$
        \item Любая подгруппа $G$, если $G$~--- Абелева.
        \item $A_n$ (знакопеременная группа степени n~--- четные перестановки)
    \end{enumerate}
\end{Example}

\begin{Lemma}{}{}
    $(G:H) = 2 \Rightarrow H \vartriangleleft G$
\end{Lemma}

\begin{proof}

    $g \in G, \ \ g \in H \Rightarrow gH = H = h'g$

    $g \not \in gH \not = H \underset{\text{классов всего 2}}{\Rightarrow} gH = G \underset{\text{разность}}{\setminus} H = Hg$
\end{proof}

\begin{Example}{Не нормальная группа}{}

    $G := S_3$ (перестановки из трех элементов)

    $H = <(1, 2)> = \{(1,2),e\}$

    $h = (1, 2), g = (1, 2, 3)$

    $ghg^{-1} = (23) \not \in H \Rightarrow H$~--- не нормальная группа в $G$.

\end{Example}
